\chapter{补充推导与实现对应}
\section{对称设定下的上下尾镜像}
本节证明命题\ref{prop:symmetry}。在本文的连续t-Copula中，联合密度存在且分位数唯一。设目标概率坐标为$V$，来源为$U$，其copula密度满足中心对称性
\begin{equation}
c(u,v)=c(1-u,1-v).
\end{equation}
由于$U$的边际密度为1，$V$在$U=\alpha$下的条件分布函数为$H_{\alpha}(v)=\int_0^v c(\alpha,w)\,dw$。作变量代换得到
\begin{align}
H_{1-\alpha}(1-v)
&=\int_0^{1-v}c(1-\alpha,w)\,dw\\
&=\int_v^1c(\alpha,z)\,dz
=1-H_{\alpha}(v).
\end{align}
因此$H_{1-\alpha}^{-1}(1-\beta)=1-H_{\alpha}^{-1}(\beta)$。目标边际关于$m_i$对称意味着$F_i^{-1}(1-v)=2m_i-F_i^{-1}(v)$，故
\begin{equation}
q_{i\leftarrow j}(1-\beta\mid1-\alpha)=2m_i-q_{i\leftarrow j}(\beta\mid\alpha).
\end{equation}
同理，当参考态为$\alpha=0.5$时，该关系依然成立。将困境态与参考态相减，得到$\Delta^+=-\Delta^-$。这一性质不要求两只股票具有相同边际自由度，只要求各自用于该条件方向的目标边际对称。上下行都使用同一预测起点的参数，因此参数随时间变化不会打破逐起点的镜像性质。

若上、下尾阈值不互为镜像，或两侧使用不同参考态，则不能套用这个等式。数值实现中的差异应先检查这些定义，而不是直接归因于市场非对称性。本稿的参考测试覆盖不同相关系数、边际自由度、条件位置和尺度；其用途是检查实现与命题的一致性，而不是替代真实市场检验。

\section{条件方向的尺度分解}
设两端条件边际为$F_i(x)=G((x-m_i)/\sigma_i)$和$F_j(x)=G((x-m_j)/\sigma_j)$，其中$G$相同。交换来源与目标时，由copula可交换性，相同$\alpha,\beta$得到同一条件概率位置$v_{\alpha,\beta}$，中位参考态位置为$v_{0.5,\beta}$。于是
\begin{align}
\Delta^-_{i\leftarrow j}
&=\sigma_i\left\{G^{-1}(v_{\alpha,\beta})-G^{-1}(v_{0.5,\beta})\right\},\\
\Delta^-_{j\leftarrow i}
&=\sigma_j\left\{G^{-1}(v_{\alpha,\beta})-G^{-1}(v_{0.5,\beta})\right\}.
\end{align}
令$k_{ij}=G^{-1}(v_{0.5,\beta})-G^{-1}(v_{\alpha,\beta})$，即得命题\ref{prop:scale}。位置参数在增量中抵消。这个抵消也说明，主边特征的底层参数控制不需要为了重建增量而额外引入条件均值；不过若改为使用CoVaR水平本身，则条件均值必须加入对应控制。

当同一目标的正强度满足$s_{i\leftarrow j}=\sigma_i k_{ij}$时，行和归一化权重为
\begin{equation}
\frac{s_{i\leftarrow j}}{\sum_{\ell}s_{i\leftarrow\ell}}
=\frac{k_{ij}}{\sum_{\ell}k_{i\ell}},
\end{equation}
因此不能同时声称“该归一化去除了个股尺度影响”与“归一化后的边仍保留了全部原始尺度信息”。类似地，top-$k$所选关系可能不互相保留；因此即使$k_{ij}=k_{ji}$，支持矩阵$A$仍可能不对称。正文的方向分析同时报告权重与支持，避免混淆。

若两端标准化边际形状不同，目标逆分位变换也不同，尺度标准化不能保证双向完全相等。本稿要求保留两端边际自由度，并进行共同边际形状控制，而不是将标准化后残留的一切差异都解释成独立方向。

\section{汇总不决定逐边条件信息：一个构造}
以下数字仅为代数示例，不是公司数据或实证结果。设节点1、2为来源，节点3、4为目标，唯一四条边的权重为
\begin{equation}
W=\begin{pmatrix} a&b\\ b&a\end{pmatrix},\qquad
W'=\begin{pmatrix} b&a\\ a&b\end{pmatrix},
\end{equation}
其中行对应目标3、4，列对应来源1、2。两个网络的每个目标入权重和、每个来源出权重和、度数及各目标权重集合均相同。然而，当来源状态为$x=(x_1,x_2)^\top$时，两个网络对目标3的线性消息差为$(a-b)(x_1-x_2)$。例如$a=0.8,b=0.2,x_1=1,x_2=4$，该目标消息分别为1.6和3.4。

因此，即使完整保留每个节点的风险总量，也可能无法恢复哪条风险关系对应哪个来源状态。该构造仅说明一般关系表示中的信息压缩，不说明这四条边能由任意给定的t-Copula参数生成，也不证明其预测价值。在真实实验中必须通过同信息集的汇总模型、等容量模型及关系重配实验验证。

\section{参考价收益、调整收益与总回报}
令$P_{t-1}$为前一交易日未复权收盘，$P_t$为当前收盘，$B_t$为当日交易所参考价。若复权价格定义为$\widetilde P_t=A_tP_t$，则其对数收益为
\begin{equation}
R_t^{\mathrm{adj}}=\log\frac{P_t}{P_{t-1}}+\log\frac{A_t}{A_{t-1}}.
\end{equation}
它等于候选参考价收益$\log(P_t/B_t)$的条件为
\begin{equation}
\frac{A_t}{A_{t-1}}=\frac{P_{t-1}}{B_t}.
\end{equation}
数据库中名为ratio的字段可能保存此比值或其倒数，累计因子也可能采用不同的锚定方向；必须从定义与样例确认，不能只凭名称代入。

仍用纯构造例子说明：若前收10元、每股现金分红1元、参考价9元且当日收盘9.9元，则参考价对数收益为$\log(9.9/9)$，若暂忽略税费和再投资，持有股票并收到1元现金的单期财富对数增长为$\log((9.9+1)/10)$。二者分别对应10\%与9\%的简单变化，因而并不相同。这不评价哪种输入对预测更好，只说明“价格调整收益”与“投资者总回报”不能未经定义直接等同。复杂送转、配股或多阶段事件更应使用明确的经济条款与生效日期。

同一会话内若所有价格统一乘以相同正因子，式\eqref{eq:intradayreturn}的差分与RV保持不变。这种不变性并不证明跨日收益正确，也不能认证整个历史复权账本的当时可得性。

\section{信息时点、标签成熟与历史归档}
设预测起点$s$对应目标区间$[s+1,s+h]$。训练截止为$c$时，标签可用至少要求$s+h\leq c$，并进一步要求目标数据已经按合同实际可用。以交易日序号表示，当$h=5,c=104$时，起点99的目标100至104已完整；起点100的目标101至105尚未成熟。判断不能依据起点100是否早于104。

历史图$G_s$由起点$s$可用的参数、状态与数据版本生成。当前$c$时点允许用已取得的历史数据重新估计当前状态，但不能将这组较新参数写回$G_s$，然后把回写后的图当成当时的输入。五日边变化必须比较$G_t$与实际归档的$G_{t-5}$；底层参数控制也从这两个归档版本取值。

证券集合在最初形成期结束时固定，早于第一条预测模型训练样本的起点。该选择只依据形成期已有的资格、流动性和覆盖信息，不根据后续完整存活或测试结果反选。若历史资格只有事后回填，数据论文必须将样本称谓与时点限制说明清楚，不能依靠“历史日期”三个字替代当时可见性证据。

\section{公式与代码的对应要求}
\begin{longtable}{p{0.30\textwidth}p{0.62\textwidth}}
\toprule 正文对象 & Codex实现及验收重点\\\midrule\endhead
式\eqref{eq:rv}、\eqref{eq:rq} & 会话锚点、完整规则采样、午休与集合竞价口径；现有参考函数只接受已整理的价格路径。\\
式\eqref{eq:target} & 用完整交易日历构造未来窗口；缺失一日不压缩窗口；保存成熟日期。\\
式\eqref{eq:unit_t} & 明确普通t与单位方差t转换，边际自由度大于2。\\
式\eqref{eq:dcc} & 正定约束、分段初始化、过去估计窗口；完整DCC估计器仍需实现。\\
式\eqref{eq:conditional_t}—\eqref{eq:quantile_back} & 从同一无序对生成双向条件量；参考测试核对镜像与尺度性质。\\
式\eqref{eq:edgefeatures} & 历史归档五日差、缺边掩码、训练尺度；底层参数对照包含同样滞后。\\
式\eqref{eq:message} & 行为目标、列为来源；编号置换等变，空邻居退回自身分支。\\
式\eqref{eq:qlike} & 正标签域、预测下限、原始与去常数损失区分；全模型共同评价样本。\\
\bottomrule
\end{longtable}
当前附带的数值代码只验证若干定义、代数性质和构造例子，未实现完整数据库读取、GARCH/DCC估计、图神经网络训练或统计实验。真实数据的验收与全部预测实验由后续研究阶段完成。
